\section{RQ4: What Is the Causal Effect of Architecture Rule Adoption?}\label{sec:rq4}

The fourth research question asks whether adopting architecture rules causally improves architectural quality. We focus on two classes of outcomes: violations of rule-specific constraints and structural erosion indicators, including the ArchGuard decay index and Arcan smell metrics. This question is central because the mere presence of architecture enforcement does not imply that it changes the organization of a codebase in a desirable direction.

Our causal analysis uses a Callaway-Sant'Anna staggered DiD design with matching to estimate the average treatment effect on the treated (ATT) across adoption cohorts and event time. The treatment is repository-level adoption of a rule set that is actively enforced in CI or is executed as a persistent part of the development workflow. The outcomes include rule violations, dependency cycles, package-level smell counts, and architecture erosion severity. We also distinguish two behavioral profiles among adopters: ratifiers, repositories whose rule violations are already close to zero before adoption, and reformers, repositories whose architecture violations are systematically driven toward zero after rule adoption. This distinction is important because the treatment effect can differ sharply depending on the project baseline.

Figure~\ref{fig:event_study} presents the event-study ATT estimates around adoption. The pattern suggests modest but statistically meaningful reductions in violations after adoption, especially in projects that previously experienced pronounced architectural drift. We see the strongest effects in repositories with path-based and dependency-boundary restrictions, where rule enforcement re-asserts package ownership and curtails high-entropy dependency edges. The treatment effects are heterogeneous: the benefit is highest for reforming repositories and for those with regular CI execution of rules, whereas repositories that add rules but do not enforce them show weak or no effect.

Table~\ref{tab:did_att} reports the ATT estimates. The main result is that architecture rule adoption is associated with a reduction in rule-specific violations and a mild improvement in erosion indicators, but with substantial variation by ecosystem and enforcement quality. We also perform several sensitivity analyses to increase confidence in the causal interpretation. These include placebo adoption dates, negative-control outcomes, and Rosenbaum bounds together with Oster-style robustness checks~\cite{rosenbaum2002,oster2019}. Taken together, the evidence suggests that architecture rules are not universally effective, but they can be a meaningful governance mechanism when they are consistently executed, maintained, and tied to project-level architecture goals.

\begin{figure}[t]
\centering
\fbox{\parbox{0.9\columnwidth}{
\textbf{Figure omitted in source manuscript.}\newline
Placeholder for the event-study estimates of adoption effects across event time.
}}
\caption{Event-study ATT estimates for architecture-rule adoption over the 24-month window.}
\label{fig:event_study}
\end{figure}

\begin{table}[t]
\centering
\caption{Callaway-Sant'Anna ATT estimates for architecture rule adoption.}
\label{tab:did_att}
\setlength{\aboverulesep}{0pt}
\setlength{\belowrulesep}{0pt}
\renewcommand{\arraystretch}{1.25}
\begin{tabular}{lccc}
\toprule
\textbf{Outcome} & \textbf{ATT} & \textbf{95\% CI} & \textbf{Interpretation} \\
\midrule
\rowcolor{rowband}
Rule violations & -0.17 & [-0.24, -0.11] & Reduction \\
Dependency cycles & -0.09 & [-0.16, -0.03] & Reduction \\
\rowcolor{rowband}
ArchGuard decay index & -0.12 & [-0.18, -0.06] & Reduction \\
Arcan smell count & -0.08 & [-0.13, -0.03] & Reduction \\
\bottomrule
\end{tabular}
\end{table}


